\section{Cost and coordination efficiency}
\label{sec:cost}

\subsection{Efficiency metrics against capacity}

Figure~\ref{fig:efficiency} shows the three coordination metrics against capacity for the
three multi-agent topologies.

\begin{figure}[H]
\centering
\includegraphics[width=\linewidth,height=0.30\textheight,keepaspectratio]{fig_efficiency.pdf}
\caption{Coordination efficiency, error amplification and turn overhead against model size,
by topology. The dashed lines mark the single-agent reference at $A_e = 1$ and
$O\% = 0$.}
\label{fig:efficiency}
\end{figure}

Coordination efficiency is 0.198 for independent, 0.102 for decentralized and 0.097 for
centralized, against 0.564 for the single-agent reference. The ordering follows the turn
count directly, since $E_c$ divides accuracy by relative turns: independent spends three
turns, both communicating topologies spend six. No multi-agent topology recovers its turn
cost in accuracy under this metric, which is the expected result given that accuracy gains
are between 0.01 and 0.08 while turn counts triple or sextuple.

Error amplification is below one for all three multi-agent topologies: 0.886 for
decentralized, 0.928 for centralized and 0.934 for independent. Every coordination
structure suppresses some fraction of the errors a single agent of the same capacity would
make, and debate suppresses the most. Against the single-agent baseline, the reduction in
error amplification from adding agents is 0.025, 0.093, 0.066 and 0.137 on GPQA, MMLU-Pro,
TruthfulQA and MATH.

Turn overhead is exactly 200 percent for independent and 500 percent for both communicating
topologies, since the counts are structural rather than data-dependent.

Moving from independent to decentralized reduces coordination efficiency by 0.067, 0.104,
0.104 and 0.112 across the four benchmarks. Moving from independent to centralized reduces
it by a similar amount while also reducing accuracy on the mathematical benchmark.

\subsection{The compute frontier}

The cost proxy is the product of parameter count and mean total generated tokens, summed
over agents and rounds. Under this proxy the four topologies cost 15{,}339, 56{,}177,
102{,}771 and 120{,}897 on average, in the order single agent, independent, centralized,
decentralized.

The cost-accuracy relationship is broadly positive but with a large dominated region: many
multi-agent cells spend substantially more compute than nearby single-agent or
independent-vote cells without reaching the empirical frontier. The atlas contains the
frontier plot and the iso-accuracy surface with multi-agent cells overlaid; the
substitution structure is in Figure~\ref{fig:equivalence-arrows} of
Section~\ref{sec:scaling}.

The mathematical benchmark is the most favourable to multi-agent systems on this frontier.
GPQA and TruthfulQA more often reward capacity and reasoning over orchestration.

\subsection{Matched-compute comparison}

The direct test is to compare three 8B agents against one 32B agent, whose parameter counts
are 24.6 billion and 32.8 billion respectively and whose costs are therefore roughly
comparable.

Across twelve contrasts spanning benchmarks and configurations, the multi-agent system is
better in 3 and worse in 7. The mean accuracy differences are $-0.121$ on GPQA, $-0.117$ on
TruthfulQA, $-0.034$ on MMLU-Pro and $+0.059$ on MATH.

On the two benchmarks where the sweep has the largest accuracy range, spending the compute
on one larger model rather than three smaller ones is better by more than ten accuracy
points. The mathematical benchmark reverses this, which is consistent with every other
result in the sweep pointing the same way: aggregation helps most where solutions are long,
verifiable and independently derivable, and least where the task is recall of a specific
fact that a larger model either knows or does not.

The comparison carries one caveat. Some of the 32B long-reasoning baselines are
context-capped, which depresses the single-agent side of the contrast in exactly the
configurations where it should be strongest. The true advantage of the larger model is
therefore likely understated. A cleaner matched-compute pair would be three 4B agents
against one 14B agent, neither of which is capped; this is recorded as a follow-up.

\subsection{Cost of the individual axes}

The matched cost contrasts quantify what each axis buys and what it costs.

Adding agents raises the cost proxy by 84{,}758, 58{,}360, 35{,}356 and 62{,}219 on the four
benchmarks while raising accuracy by 0.014, 0.031, 0.019 and 0.056.

Moving from independent to decentralized approximately doubles the cost, since it doubles
the round count, and buys between 0.007 and 0.016 accuracy.

The reasoning axis buys 0.087 to 0.135 accuracy at b2048 against disabled thinking. Since
thinking tokens dominate total tokens at the larger rungs, the cost multiplier is
substantial, but the axis is still favourable relative to orchestration on three of four
benchmarks: b2048 gives roughly five to nine times the accuracy gain of moving from
independent to decentralized.

The prompt axis buys 0.015 to 0.070 accuracy for a few hundred additional prompt tokens and
no additional generation turns. On a cost-normalised basis it is the most favourable
intervention in the sweep by a wide margin, and it is the only axis that improves accuracy
and agreement calibration simultaneously on every benchmark.

Appendix~\ref{app:marginals} contains the matched cost contrasts for every axis and the
marginal cost, total-token and reasoning-token means for every level.
